\Needspace{6\baselineskip}
\begin{center}\hypertarget{native-265-article191}{}\pdfbookmark[2]{191.}{article-191}191.\end{center}
Developing the theory of perturbations from given elements falls outside
our purpose. Here we need only show how to correct an approximate orbit so
that, when combined with given perturbations, it fits the observations as
closely as possible. The simplest way to do so is by a method analogous to
those we explained in Arts.~124, 165, 188. For the times of all the observations
we intend to use for this purpose, which may number three, four, or more
as the circumstances require, calculate the numerical values of the
perturbations from their equations, for the longitudes in the orbit, the
radius vectors, and the heliocentric latitudes. Take the arguments for this
calculation from the approximate elliptic elements on which the theory of
perturbations was based. Then select two of the observations and calculate
the distances from the Earth for them from the same approximate elements.
These distances form the first hypothesis. The second and third hypotheses
are formed by slightly changing the distances. Under each hypothesis,
determine the heliocentric positions and the distances from the Sun from
the two geocentric positions. Once the perturbations have been removed from
the latitudes, use the heliocentric positions to derive the longitude of the
ascending node, the inclination of the orbit, and the longitudes in the orbit.
In this calculation the method of Art.~110 requires some modification if
it seems worthwhile to take account of the secular variation of the
longitude of the node and of the inclination. Let $\beta,\beta'$ denote
the heliocentric latitudes with the periodic perturbations removed,
$\lambda,\lambda'$ the heliocentric longitudes, $\AscendingNode,\AscendingNode+\Delta$
the longitudes of the ascending node, and $i,i+\delta$ the inclinations of
the orbit. It is convenient to write the equations in the following form.
\[
\begin{gathered}
\mathrm{tang}\,\beta=\mathrm{tang}\,i\sin(\lambda-\AscendingNode)\\
\frac{\mathrm{tang}\,i}{\mathpunct{.}\mathrm{tang}(i+\delta)}\mathrm{tang}\,\beta'
=\mathrm{tang}\,i\sin(\lambda'-\Delta-\AscendingNode).
\end{gathered}
\]
\SourceLeaf{266}Here the value of $\frac{\mathrm{tang}\,i}{\mathrm{tang}(i+\delta)}$
is obtained to all the precision required by substituting an approximate
value for $i$. Then $i$ and $\AscendingNode$ can be found by the usual methods.

Next subtract the sums of the perturbations from the two longitudes in
the orbit and from the two radius vectors to obtain the purely elliptic
values. The effect of the secular variations in the position of perihelion
and in the eccentricity on the longitude in the orbit and the radius vector
is to be determined by the differential formulas of Section~I of Book One.
It must be included with the periodic perturbations at once if the
observations are far enough apart for this effect to be worth taking into
account. From these corrected longitudes in the orbit and radius vectors,
together with the corresponding times, determine the remaining elements.
Finally, calculate the geocentric positions for all the remaining
observations from these elements. Compare them with the observed positions
in the manner explained in Art.~188 to find the set of distances that yields
elements fitting all the remaining observations as closely as possible.
